\documentclass{article}
\usepackage[margin=28mm]{geometry}
\usepackage{amsmath,amssymb}
\usepackage{xurl}
\usepackage[hidelinks]{hyperref}
% Shared commands for notes in docs/paper-gaps/.

\DeclareUnicodeCharacter{2080}{\ifmmode{}_{0}\else\textsubscript{0}\fi}
\DeclareUnicodeCharacter{2081}{\ifmmode{}_{1}\else\textsubscript{1}\fi}
\DeclareUnicodeCharacter{2082}{\ifmmode{}_{2}\else\textsubscript{2}\fi}
\DeclareUnicodeCharacter{2099}{\ifmmode{}_{n}\else\textsubscript{n}\fi}

\newcommand{\Lean}{\textsc{Lean}}
\ExplSyntaxOn
\NewDocumentCommand{\leanid}{m}{
  \ifmmode
    \textsf{\detokenize{#1}}
  \else
    \group_begin:
    \urlstyle{sf}
    \tl_set:Nn \l_tmpa_tl {#1}
    \tl_replace_all:Nnn \l_tmpa_tl {\_} {_}
    \exp_args:NV \path \l_tmpa_tl
    \group_end:
  \fi
}
\ExplSyntaxOff
\providecommand{\tr}{\operatorname{tr}}
\newcommand{\tnleanIssueBase}{https://github.com/LionSR/TNLean/issues/}
\newcommand{\tnleanPrBase}{https://github.com/LionSR/TNLean/pull/}
\newcommand{\tnleanDocBase}{https://sirui-lu.com/TNLean/docs/}
\newcommand{\ghissue}[1]{\href{\tnleanIssueBase#1}{GitHub issue \##1}}
\newcommand{\ghpr}[1]{\href{\tnleanPrBase#1}{PR \##1}}
\newcommand{\doclink}[1]{\href{\tnleanDocBase#1.html}{\path{#1}}}

% Dirac notation and the identity operator, as in blueprint/src/macros/common.tex.
% \providecommand keeps note-local definitions authoritative.
\usepackage{dsfont}
\providecommand{\ket}[1]{|#1\rangle}
\providecommand{\bra}[1]{\langle #1|}
\providecommand{\braket}[2]{\langle #1 | #2 \rangle}
\providecommand{\ketbra}[2]{|#1\rangle\!\langle #2|}
\providecommand{\Id}{\mathds{1}}
\providecommand{\one}{\mathds{1}}
\providecommand{\C}{\mathbb{C}}
\providecommand{\MN}[1]{M_{#1}(\C)}
\providecommand{\MD}{M_D(\C)}

% Verdict marker: \gapnote{<kind>}{<status>}. Kind is the policy
% classification (clarification, local-correction, scope-restriction,
% unfaithful, false-source, open-gap); status is open, wip (elimination
% underway), resolved, or historical. Machine-read by the site index;
% prints a verdict line.
\newcommand{\gapnote}[2]{\par\noindent\textbf{Verdict:} #1 (#2).\par}

\providecommand{\gapnote}[2]{\par\noindent\textbf{Verdict:} #1 (#2).\par}
\title{Cluster Size and Diameter in Hastings' Percolation Bound}
\author{}
\date{}
\begin{document}
\maketitle
\gapnote{local-correction}{resolved}

\section{The printed step}
Hastings, \emph{Solving gapped Hamiltonians locally}~\cite{Hastings2006Solving},
\href{https://arxiv.org/abs/cond-mat/0508554v3}{arXiv:cond-mat/0508554v3},
source \path{proj3new.tex}, lines 1227--1254, bounds the sum over clusters of
diameter larger than $l_{\rm proj}$ by
\begin{align}
 \sum_{m\ge1}\frac1{m!}\Bigl(\sum_{\rm clusters}y^{n_C}\Bigr)^m
 \le\exp\Bigl[N\sum_{n=l_{\rm proj}+1}^{\infty}(xy)^n\Bigr]-1,
 \label{eq:hastings05_percolation_printed}
\end{align}
where $n_C$ is the number of sites of the cluster $C$ and at most $Nx^n$
clusters have $n$ sites (line 1240). The sum starts at $n=l_{\rm proj}+1$,
which presumes that a cluster of diameter larger than $l_{\rm proj}$ has more
than $l_{\rm proj}$ sites.

\section{The relation between size and diameter}
A cluster is a set $C$ of sites $i$ whose terms $H_i$ are joined by
overlapping supports, and $H_i$ acts on sites within distance $R$ of $i$
(lines 1131--1139). Consecutive members of a cluster can therefore be $2R$
apart. If the cluster has $n$ sites, the sites acted on by its terms span a
distance of at most $2Rn$, since each support has diameter at most $2R$ and
the supports form a connected chain. Hence a cluster of diameter larger than
$l_{\rm proj}$ satisfies
\begin{align}
 n_C\ge\Bigl\lfloor\frac{l_{\rm proj}}{2R}\Bigr\rfloor+1,
 \label{eq:hastings05_percolation_size}
\end{align}
and no more holds in general: on a chain with $R\ge1$ the cluster $\{0,2R\}$
has $n_C=2$, since $H_0$ and $H_{2R}$ share the site $R$, yet it spans the
distance $4R$, which exceeds $l_{\rm proj}=4R-1\ge3$. The printed start
$n=l_{\rm proj}+1$ is right for nearest-neighbour bonds, as in the figure at
lines 1120--1129, and not for $R>1$.

\section{The correction}
Start the sum at $n=l+1$ with $l=\lfloor l_{\rm proj}/(2R)\rfloor$, as
(\ref{eq:hastings05_percolation_size}) allows. For $0<xy\le1/2$,
\begin{align}
 N\sum_{n\ge l+1}(xy)^n\le N(xy)^l
 \le N(xy)^{-1}e^{-l_{\rm proj}/\xi_{\rm perc}},
 \qquad
 \xi_{\rm perc}=\frac{2R}{\ln(1/(xy))},
\end{align}
since $l\ge l_{\rm proj}/(2R)-1$. The decay length is of order the interaction
range, as the source states at line 1250; the factor $R$ there is what
compensates the printed start of the sum. Only the constant $N$ changes, to
$N(xy)^{-1}$.

\section{Scope}
The formalized percolation sum takes the size threshold $l$ as a hypothesis
on every cluster and states its bound as $\exp(Ne^{-l/\xi})-1$ with
$\xi=1/\ln(1/(xy))$, a decay length in units of cluster size. The
conversion to distance is the inequality
(\ref{eq:hastings05_percolation_size}).
\bibliographystyle{unsrt}
\bibliography{references}
\end{document}
